\chapter{Conclusion and Perspectives}

This thesis investigated the use of logic programming as an implementation
framework for elaboration mechanisms in \rocq, with a particular focus on
type-class resolution and static analyses. More broadly, our goal was to
explore how an external logic-programming engine can be used both as an
implementation tool and as a framework for reasoning about components of an
interactive theorem prover.

A first contribution of this work is the implementation of an alternative
type-class solver in \elpi. We showed that type classes and instances can be
compiled into a logic-programming database using a relatively small amount of
code, while relying on the existing \elpi interpreter for resolution. The
resulting solver shows competitive performance with respect to the native
\rocq implementation. Experimental results presented in
\cref{sec:bench} indicate that the overhead introduced by translating terms
between the object language and the meta-language remains limited in practice,
making this approach viable from an efficiency perspective.

The integration of \rocq and \elpi is, however, not without challenges.
Although both systems support higher-order unification and
$\beta\eta$-conversion, their internal representations of terms differ
significantly. This mismatch leads to situations where terms that unify in
\rocq fail to unify after translation into \elpi. To address this issue, we
developed a compilation scheme based on links and delayed constraints. This
approach allows us to recover much of the intended behavior of \rocq
unification while preserving the advantages of the \elpi execution model,
including indexing and mode-directed unification.

Beyond type-class resolution, this thesis also investigated static analyses for
logic programs. We introduced a determinacy analysis for \elpi capable of
handling higher-order predicates, dynamic clauses, and the cut operator. The
analysis provides guarantees that predicates declared deterministic do not
introduce unnecessary choice points during execution. In the context of
type-class resolution, such guarantees contribute to a better understanding and
control of the search performed by the logic-programming engine.

Finally, we initiated the formal verification of these results by mechanizing a
first-order fragment of the \elpi interpreter in \rocq. Although this model
does not capture the full higher-order language, it is sufficiently expressive
to establish machine-checked correctness results for the determinacy analysis.
This mechanization constitutes a first step towards a more comprehensive formal
account of \elpi and its analyses.

Taken together, these results suggest that logic programming is not merely a
convenient implementation technique for proof assistants, but also a useful
framework for studying, prototyping, and formally reasoning about elaboration
mechanisms. The interaction between \rocq and \elpi demonstrates that existing
logic-programming technology can be leveraged effectively while still preserving
the level of rigor expected from theorem-proving environments.
